\cotefile{93}
\scanpage{47}
peut énoncer qu'il y a un isomorphisme canonique (fonctoriel, compatible
aux extensions de la base)
\[
\deleted{\unsure{\underline E(}}\handwritten{\mathbb D}^{*}(A)=\underline{\mathrm{Lie}}(\underline E(A^{*}))\ ,
\]
la structure d'extension de $\mathbb D^{*}(A)$ n'étant autre que la structure\inserted{\handwritten{sur \uline{Lie}}}
déduite de la structure d'extension de $\underline E(A)_S=E(A)$.

Ceci posé, revenant aux conditions du théo\deleted{\unsure{èr|rè}}rème de déformations, on
voit, par un sorite sur l'exponentielle sur lequel nous reviendrons, que
la donnée d'un prolongement $\mathrm{Fil}^{1}\handwritten{\subset}\mathbb D^{*}(A)_{S'}=\underline{\mathrm{Lie}}(\underline E(A^{*})_{S'})$ localement facteur
direct prolongeant $\underline{\mathrm{Lie}}(\underline{\underline{t}}^{\handwritten{\vee}}_A)=\underline t^{\handwritten{\vee}}_A\handwritten{\subset}\mathbb D^{*}(A)_S=\underline{\mathrm{Lie}}(\underline E(A^{*})_{S'}\handwritten{\mid} S)$ équivaut à
la donnée d'un sous-groupe \inserted{lisse} a structure vectorielle de $\underline E(A^{*})_{S'}$ qui prolonge le sous-groupe $\underline t_A^{\handwritten{\vee}}$ de $\underline E(A^{*})_S$. Si $L$ est ce sous-groupe, alors il est
immédiat que \deleted{$\unsure{A^{*}(|E^{*}(}$}
\[
\underline E(A^{*})_{\deleted{\unsure{t}}\handwritten{S'}}/L
\]
est un schéma abélien sur $S'$, dont le schéma abélien dual est le $A'$ cherché. (NB Il serait plus naturel ici de travailler avec \deleted{$\unsure{D_{*}(A)=D(A)|\mathbb D^{*}(A)=\mathbb D(A)}$}
\[
\mathbb D_{*}(A)=\underline{\mathrm{Lie}}\ \underline E(A)\ \overset{\text{isom can}}{\handwritten{\xrightarrow{\sim}}}\ \mathbb D^{*}(A^{*})^{\handwritten{\vee}}
\]
plutôt qu'avec $D^{*}(A)$ pour énoncer le théorème de déformations; l'équivalence des deux points de vue provient de l'accouplement parfait
\[
\mathbb D^{*}(A)\otimes\mathbb D^{*}(A^{*})\handwritten{\longrightarrow}\underline O_{S_{\mathrm{c}\deleted{\unsure{\mathrm{ii}}}\mathrm{risnilp}}}
\]
compatible avec les filtrations.\handwritten{)}

\uline{Remarque}. En t\deleted{\unsure{éh}}héorie des groupes de BT, nous donnerons deux constructions de $\handwritten{\mathbb D}^{\handwritten{*}}(G)$, inspirées des deux constructions indiquées dans le cas des
schémas abéliens: l'une "cohomologique", l'autre via une "extension vectorielle universelle" $\underline E(G)$. La construction de $\underline E(G)$ peut aussi s'obtenir
par deux méthodes, l'une cohomologique, l'autre"par l'exponentielle" étant
directement adaptée à la démonstration du théorème de déformations.
% archive sheet number: 46

